% ============================================================================
\section{跳转到用户态：\texttt{iret}}
\label{sec:tss-iret}
% ============================================================================

目标：从 Ring 0 跳到 Ring 3 执行一段用户态代码。
关键是构造一个 \texttt{iret} 帧——让 CPU 以为自己在从中断返回到 Ring 3。

% ----------------------------------------------------------------------------
\subsection{\texttt{iret} 帧的构造}
% ----------------------------------------------------------------------------

\texttt{iret} 从栈上依次弹出 EIP → CS → EFLAGS。当检测到目标 CS 的 RPL > 当前 CPL 时（跨特权级返回），还会额外弹出 ESP → SS。

所以我们需要从高到低依次 push 5 个值：

\begin{figure}[H]
  \centering
  \begin{tikzpicture}[
    cell/.style={draw, minimum height=0.8cm, minimum width=5cm, anchor=west, font=\ttfamily\small},
    addr/.style={font=\ttfamily\footnotesize, anchor=east},
  ]
    \node[font=\bfseries] at (2.5, 0.5) {iret 帧（栈上布局）};

    % Stack grows downward, so low address is at bottom
    \node[cell, fill=green!15] (ss) at (0, -0.3)  {SS \quad\quad (0x23)};
    \node[cell, fill=green!10] (esp) at (0, -1.1)  {ESP \quad\quad (user\_stack\_top)};
    \node[cell, fill=yellow!15] (efl) at (0, -1.9) {EFLAGS \quad (IF=1)};
    \node[cell, fill=blue!10] (cs) at (0, -2.7)   {CS \quad\quad (0x1B)};
    \node[cell, fill=blue!15] (eip) at (0, -3.5)  {EIP \quad\quad (user\_entry)};

    % Push order annotations
    \node[addr] at (-0.2, -0.3) {先 push};
    \node[addr] at (-0.2, -3.5) {后 push};

    % Pop order
    \draw[->, thick, red!60!black] (5.5, -3.5) -- (5.5, -0.3)
      node[right, pos=0.5, text width=2.5cm, font=\small] {iret 弹出\\顺序};

    % Address annotation
    \node[font=\small, gray, anchor=west] at (5.8, -0.3) {高地址};
    \node[font=\small, gray, anchor=west] at (5.8, -3.5) {低地址 (ESP)};

    % Cross-privilege note
    \draw[decorate, decoration={brace, amplitude=5pt}, thick]
      (-0.3, 0.0) -- (-0.3, -1.4);
    \node[font=\small, text width=3cm, anchor=east, align=right] at (-0.5, -0.7)
      {跨特权级时\\额外弹出};
  \end{tikzpicture}
  \caption{\texttt{iret} 帧布局——跨特权级返回时的 5 个栈元素}
  \label{fig:iret-frame}
\end{figure}


% ----------------------------------------------------------------------------
\subsection{汇编实现：\texttt{jump\_to\_ring3}}
% ----------------------------------------------------------------------------

\codefile{src/kernel/arch/tss\_flush.asm}
\begin{lstlisting}[style=nasmstyle]
USER_CODE_SEL equ 0x1B     ; GDT[3] | RPL=3
USER_DATA_SEL equ 0x23     ; GDT[4] | RPL=3

; void jump_to_ring3(uint32_t eip, uint32_t esp3);
jump_to_ring3:
    mov eax, [esp + 4]     ; eax = eip（用户态入口）
    mov ecx, [esp + 8]     ; ecx = esp3（用户态栈顶）

    ; 先切换 DS/ES/FS/GS 到用户数据段
    mov dx, USER_DATA_SEL
    mov ds, dx
    mov es, dx
    mov fs, dx
    mov gs, dx

    ; 构造 iret 帧
    push dword USER_DATA_SEL  ; SS
    push ecx                   ; ESP
    pushfd                     ; EFLAGS
    or dword [esp], 0x200      ; 确保 IF=1
    push dword USER_CODE_SEL  ; CS
    push eax                   ; EIP

    iret                       ; 跳转到 Ring 3！
\end{lstlisting}

\begin{warnbox}
\textbf{为什么要在 \texttt{iret} 前切换 DS/ES/FS/GS？}

如果不切换，\texttt{iret} 后 DS 仍指向内核数据段（DPL=0），
但 CPL 已变为 3。用户态代码第一次访问 DS 时，CPU 检查发现
$\max(\text{CPL}=3, \text{RPL}) > \text{DPL}=0$，立即触发 \#GP。

所以必须在 \texttt{iret} 之前把 DS/ES/FS/GS 切换到 DPL=3 的
用户数据段（\hex{23}）。
\end{warnbox}

\begin{cpustate}
\texttt{iret} 执行后：
\begin{itemize}
  \item CPL 从 0 变为 3（用户态）
  \item CS 从 \hex{08} 变为 \hex{1B}（User Code）
  \item SS 从 \hex{10} 变为 \hex{23}（User Data）
  \item ESP 切换到用户态栈（参数 \texttt{esp3}）
  \item EIP 指向用户态代码入口（参数 \texttt{eip}）
  \item EFLAGS.IF = 1（中断保持打开，键盘和定时器仍然工作）
  \item 后续中断发生时，CPU 自动从 TSS 读取 SS0:ESP0 切回内核栈
\end{itemize}
\end{cpustate}

% ----------------------------------------------------------------------------
\subsection{为什么要确保 IF=1？}
% ----------------------------------------------------------------------------

\texttt{pushfd} 保存的 EFLAGS 可能 IF=0（如果在关中断的上下文中调用）。
\texttt{or dword [esp], 0x200} 强制设置 bit 9（IF），确保用户态运行时
中断处于打开状态——否则键盘和定时器中断都不会触发，系统假死。

\begin{thinkbox}
如果忘记设置 IF=1，用户态代码执行了怎样的循环会导致系统完全无响应？

\emph{提示：任何循环都会，因为 IF=0 时所有可屏蔽中断（包括键盘 IRQ1）
都被禁止，没有任何机制能打断 Ring 3 的循环回到内核。}
\end{thinkbox}
